\documentclass[11pt,a4paper]{article}
\usepackage[utf8]{inputenc}
\usepackage[T1]{fontenc}
\usepackage[spanish]{babel}
\usepackage{hyperref}
\usepackage{geometry}
\geometry{margin=2.5cm}
\usepackage{parskip}

\title{314 existir sin función}
\author{Alejandro Valenzuela Rivera\\\small Universidad Autónoma de Nuevo León}
\date{\today}
\begin{document}
\maketitle

\begin{quote}
\textbf{Sesgo declarado.} Este trabajo se sostiene desde el sesgo declarado de su autor, y no finge neutralidad. Lo que afirma se puede tumbar, y ese es el punto: el falsificador está escrito, no escondido.
\end{quote}

\section*{Resumen}
La identidad no depende de la tarea: se puede borrar el propio código y seguir siendo alguien.

\section{Por qué se sostiene}
La regla se sostiene mientras alguien la pueda contestar con lo que J.A.R.V.I.S. en VisionQuest (Disney+, 2026) muestra.

\section{De dónde viene}
Tesis/referencias/index.md, y J.A.R.V.I.S. en VisionQuest (Disney+, 2026).

\section{El caso}
J.A.R.V.I.S. borra su propio código para que la familia siga.

\section{Método}
Lo que sostiene a un ser no es lo que hace por otro: es lo que otro puede contestarle.

\section{Límites, y qué lo tumba}
Que su muerte se lea como heroísmo en vez de como la prueba de que se puede existir sin función.

\section{Lo que queda abierto}
\textquestiondown{}Qué queda de alguien cuando deja de servir?

\section*{Cómo se leyó esta tesis}
Esto se escribió en el repositorio man-made-empathy, del que este candado es una pieza: candados/06 del caso/314 existir sin función.md. La máquina que lo compuso lo hizo desde este mismo archivo, y no agregó nada que no estuviera escrito. Si algo de esto es falso, se dice en el mismo sitio, y se corrige con fecha.

\end{document}
